\chapter[Data and Feedback: The Fifth Element]{Data and Feedback: The Fifth Element}
\chaptermark{The Fifth Element}
\label{ch:data}

Four elements so far: the model, the compute, the semiconductor, the watt. The rent-migration logic of the preceding chapters demands a fifth, and this chapter supplies it---closing a gap the argument would otherwise step around rather than confront. If open weights commoditize the model layer, and value migrates to adjacent scarcities (Section~\ref{sec:rent-migration}), then the most durable adjacent scarcity is the one input that openness cannot replicate by construction: \emph{data that only deployment generates}. This chapter maps what open weights do and do not commoditize on the data axis, who holds which class of scarce data, and the open questions that will decide whether feedback---not weights, not FLOPs---turns out to be the true scarce layer of the mature industry.

\section{What open weights do not commoditize}
\label{sec:data-classes}

An open-weight release transfers a compressed function of its training distribution. It does not transfer, and cannot: (1) \emph{proprietary enterprise workflows}---the tickets, codebases, contracts, and process traces inside firms, reachable only through deployment relationships; (2) \emph{scientific-instrument data}---beamline, sequencer, telescope, and sensor streams, gated by institutions and provenance; (3) \emph{clinical and laboratory data}, gated by regulation and consent; (4) \emph{high-quality human correction and preference data}, produced only where skilled users interact with deployed systems at scale; (5) \emph{embodied physical-interaction data}, produced only by deployed fleets (Chapter~\ref{ch:embodied}); (6) \emph{evaluation and failure traces}---what actually broke, in production, under adversarial pressure; and (7) \emph{operational telemetry} from large deployed systems. Each is generated by \emph{use}, appropriated by whoever owns the deployment relationship, and---critically for this book---each remains scarce \emph{after} the model becomes free. Table~\ref{tab:data-classes} assigns the current holders.

\begin{table}[htbp]
  \centering
  \small
  \caption{Data classes after model commoditization: scarcity mechanism and current holder of advantage [A/S].}
  \label{tab:data-classes}
  \setlength{\tabcolsep}{4pt}
  \begin{tabular}{p{36mm}p{48mm}p{48mm}}
    \toprule
    Data class & Scarcity mechanism & Likely advantage \\
    \midrule
    Public pretraining corpora & increasingly replicated; approaching exhaustion & the commons---no one's rent \\
    \addlinespace
    Synthetic reasoning data & teacher capability and filtering quality & frontier laboratories (both blocs) \\
    \addlinespace
    Enterprise workflow data & access and integration; deployment relationships & US cloud and application firms \\
    \addlinespace
    Scientific and instrument data & provenance; public institutions; consent & national laboratories and consortia \\
    \addlinespace
    Industrial and embodied data & manufacturing deployment surface & China (Chapter~\ref{ch:embodied}) \\
    \addlinespace
    Human preference and correction data & skilled-user deployment at scale & distribution owners (currently US consumer, mixed enterprise) \\
    \addlinespace
    Evaluation and incident data & certification and disclosure regimes & \emph{currently unowned}---the same vacancy as Chapter~\ref{ch:trust}'s certification layer \\
    \bottomrule
  \end{tabular}
\end{table}

Read as a campaign map, the table explains several things the four-element story left dangling. It explains why the American flywheel's rent-recapture plan (Table~\ref{tab:rent}) can survive open weights: the enterprise-workflow row is the one Chinese open models cannot reach from outside the deployment relationship, and it is the row Anthropic's \$47B enterprise run rate is built on. It explains why China's embodied push (Chapter~\ref{ch:embodied}) is strategically rational even if humanoids disappoint commercially: the industrial-data row is the only row where China holds a structural generation advantage, and the state data-farm program is an attempt to \emph{manufacture} a data moat the way Part~I manufactured cost moats. And it explains the last row's importance one more time: evaluation and incident data---who failed, how, under what attack---is the raw material of the certification authority this book keeps finding unowned.

\section{The synthetic-data question}
\label{sec:synthetic}

The sharpest challenge to any data-moat argument is synthetic data: if strong models can generate their own training material, deployment-derived data loses scarcity value. The evidence cuts both ways, and the honest statement is conditional. Synthetic reasoning data demonstrably works---the distillation economy of Chapter~\ref{ch:model} is built on it, post-training pipelines are majority-synthetic at every major lab, and the alleged industrial-scale harvesting of frontier outputs (Chapter~\ref{ch:trust}) is precisely a fight over synthetic-data \emph{feedstock}. But synthetic data has a teacher: its quality is bounded by the capability of the model that generates it plus the selectivity of the filter that curates it, which makes it a mechanism for \emph{diffusing} frontier capability rather than exceeding it---distillation relocates the frontier's rent, it does not abolish the frontier. And recursive training on model outputs degrades distributional diversity and factual grounding unless anchored by fresh non-synthetic signal; the model-collapse literature overstates the cliff, but the direction is real. The synthesis this book adopts [S]: synthetic data \emph{amplifies} whoever already holds either frontier capability (the teacher) or verified ground truth (the anchor)---it does not substitute for the anchor. Which returns the argument to the table above, because the anchors are exactly the deployment-generated classes: verified task outcomes, experimental results, physical interaction, production failures. \emph{Verification, not generation, is the scarce half of the synthetic-data economy}---a conclusion that should look familiar, because it is the funnel of Chapter~\ref{ch:prize} ($v_j$ before $s_j$) and the trust stack of Chapter~\ref{ch:trust} restated as a data-supply question.

\section{Seven questions that decide the element}
\label{sec:data-questions}

The element poses seven open questions, and this book puts each in the sharpest form it can before answering it to the extent the evidence permits, registering the rest.

\emph{Who owns data generated by deployed agents?} Legally unsettled everywhere, and the default is being set by contract while legislatures sleep: enterprise agreements increasingly reserve workflow traces to the customer while permitting aggregate learning to the provider---which, at scale, gives the provider the moat and the customer a receipt. The jurisdiction that defines agent-generated data rights first will set the default for everyone, and none has.

\emph{Can open models be improved with closed workflow traces without leaking proprietary knowledge?} Technically yes---LoRA-class adapters, federated post-training, and gradient aggregation exist---but the assurance problem (proving the adapter does not memorize the traces) is unsolved at certification grade, which makes this another product the trust layer of Chapter~\ref{ch:trust} must ship before regulated industries participate.

\emph{How should data withdrawal propagate through derivative chains?} Today: not at all. A revoked dataset persists in every model already trained on it and every derivative thereafter---the same unowned-lineage problem as the security chapter's artifact chain, and solvable by the same instrument: signed manifests and data bills-of-materials in the certification ladder.

\emph{Can scientific consortia create a data commons without making sensitive science globally downloadable?} This is the consortium design problem of Chapter~\ref{ch:consortium}, and the layered-access architecture there (open weights, tiered data, sovereign enclaves for dual-use domains) is this book's answer; the fifth element supplies its economic justification, because pooled instrument data is the one asset that makes a multi-state consortium more than the sum of its national programs.

\emph{Does industrial deployment create a physical-data advantage analogous to the US enterprise-software data advantage?} Chapter~\ref{ch:embodied}'s answer: structurally yes---same mechanism, different substrate---but only if the task-policy layer can consume it, which is the contested middle of the embodied stack.

\emph{Does synthetic data reduce or amplify frontier dependence, and when does recursion degrade?} Answered above: it amplifies whoever holds the teacher or the anchor; recursion degrades absent fresh verified signal.

\emph{Is verified feedback the true scarce layer after weights commoditize?} The book's answer, now stated as its position [S]: yes, with one refinement---the scarce object is not feedback in bulk but \emph{feedback attached to verified outcomes}, because unverified feedback is synthetic data's cheap cousin and the world is about to drown in it. Every prior element of this book converges on that refinement: the funnel's validation term, the trust stack's behavioral layer, the science metric's replication vector, the embodied proof suite's intervention-free hours. The fifth element is, in the end, the same element measured everywhere else: \emph{ground truth, at scale, with provenance}. That is what cannot be downloaded.

The dashboard rows follow directly: enterprise-agreement terms on agent-trace ownership (watch the default settle); first certification-grade privacy-preserving fine-tuning product; first data-BoM requirement in any jurisdiction's AI procurement; and the ratio of verified-outcome data to raw interaction data in published post-training recipes---the number that will reveal, before any market does, whether the fifth element is consolidating into a moat and whose.
